\chapter*{Abstract}
\addcontentsline{toc}{chapter}{Abstract}

\begin{abstract}
This interim report presents the progress made during the first semester of
a project investigating enhancement-mode operation in GaAs/AlGaAs
High-Electron-Mobility Transistors through controlled gate recession. 

The initial phase of the work focused on establishing a robust experimental and
theoretical foundation. A complete fabrication of baseline depletion-mode
HEMT devices was carried out, producing a structured array of devices with varying 
lateral geometries. These devices were electrically characterised to extract key parameters including threshold
voltage, transconductance, and gate leakage behaviour, providing a validated
reference dataset for subsequent comparison.

The results demonstrate consistent depletion-mode operation across all
fabricated devices, with threshold voltage shown to be independent of lateral
geometry and governed primarily by the vertical heterostructure beneath the
Schottky gate. This confirms gate recess depth as an appropriate process
variable for threshold-voltage engineering. Alongside experimental work,
relevant device physics and fabrication literature were reviewed to inform
process selection and experimental design.

Building on these findings, the second semester of the project will focus on
the development and calibration of a controlled gate recession process to
systematically vary donor layer thickness beneath the gate. A dedicated
device mask and fabrication strategy have been designed to enable an
empirical investigation of the relationship between recess depth and threshold
voltage, with the ultimate aim of achieving repeatable enhancement-mode
operation.
\end{abstract}
